\chapter{Discussão}

\section{Comparação entre abordagens}

\subsection{SimCLR vs CNN-JEPA vs DGAE}

As três técnicas implementadas apresentam características distintas que as tornam adequadas para diferentes cenários, cada uma com seus próprios trade-offs entre simplicidade, eficiência e qualidade de representações.

\textbf{SimCLR} apresenta a vantagem da simplicidade: a implementação é direta e o treinamento é relativamente rápido, conforme observado por \cite{chen2020simclr}. Os resultados demonstram boa separação de classes e robustez a variações através de aumentação contrastiva. No entanto, SimCLR requer aumentação explícita, o que pode ser uma limitação quando aumentação tradicional não é adequada ou quando se deseja evitar dependência de transformações específicas. Além disso, SimCLR pode ser sensível a hiperparâmetros como a temperatura do softmax, requerendo ajuste cuidadoso.

\textbf{CNN-JEPA} oferece uma alternativa interessante ao dispensar aumentação explícita, conforme proposto por \cite{kalapos2025cnnjepa}. Os resultados demonstram que é possível aprender boas representações sem aumentação tradicional, validando a abordagem de predição de embeddings conjuntos. A arquitetura é mais complexa que SimCLR, mas ainda relativamente simples comparada a outras técnicas. Uma possível desvantagem é que CNN-JEPA pode requerer mais épocas para convergência, especialmente em datasets pequenos.

\textbf{DGAE} destaca-se pela eficiência: o espaço latente é 2x menor que autoencoders tradicionais, conforme reportado por \cite{liu2025dgae}, mantendo qualidade competitiva para clustering. O guiamento por modelo de difusão permite que o autoencoder aprenda representações mais compactas e expressivas. No entanto, DGAE requer um modelo de difusão pré-treinado, o que adiciona complexidade ao pipeline. Além disso, o treinamento é mais complexo devido à integração do modelo de difusão no processo de aprendizado.

\subsection{Fatores de Decisão}

A escolha entre técnicas deve considerar múltiplos fatores, incluindo tamanho do dataset, recursos computacionais disponíveis, objetivo final da aplicação, e disponibilidade de modelos pré-treinados. Para datasets grandes com recursos computacionais abundantes, SimCLR pode ser a escolha mais simples e eficaz. Para aplicações onde aumentação tradicional não é desejável, CNN-JEPA oferece uma alternativa viável. Para aplicações onde compactação de representações é crítica, DGAE é a escolha mais adequada.

\section{Impacto da aumentação por difusão}

\subsection{Hipóteses e Validação}

Avaliamos a hipótese de que aumentação por difusão melhora a qualidade das representações aprendidas, especialmente em cenários com poucas amostras por classe. Os resultados apresentados na Tabela \ref{tab:clustering_augmented} validam parcialmente esta hipótese: a aumentação por difusão melhora consistentemente as métricas de clustering, com ganhos mais pronunciados para SimCLR.

A hipótese de que o impacto é maior em cenários com poucas amostras também foi validada: quando há poucas amostras por classe (< 50), a diversidade adicional fornecida pelas imagens geradas é mais valiosa, resultando em melhorias mais significativas nas métricas. Para datasets maiores, o ganho é menor mas ainda positivo, indicando que aumentação por difusão pode ser benéfica mesmo quando há amostras suficientes.

A hipótese de que diferentes técnicas respondem de forma diferente à aumentação foi parcialmente validada: SimCLR, que depende de aumentação tradicional, beneficia-se mais da aumentação por difusão. CNN-JEPA, que não depende de aumentação explícita, também se beneficia, mas em menor grau. DGAE, que já utiliza conhecimento do modelo de difusão, apresenta ganhos menores, como esperado.

\subsection{Análise dos Resultados}

Os resultados demonstram que aumentação por difusão pode ser uma ferramenta valiosa para melhorar qualidade de representações, especialmente quando há poucas amostras. As imagens geradas mantêm características semânticas das classes originais enquanto introduzem variação natural, propriedade que é especialmente valiosa quando aumentação tradicional pode alterar características relacionadas à classe.

No entanto, a qualidade das gerações é crucial: se o modelo de difusão não captura bem as características das classes, as imagens geradas podem introduzir ruído ao invés de diversidade útil. Portanto, é importante garantir que o modelo de difusão seja bem treinado antes de utilizá-lo para aumentação.

\subsection{Interpretação}

A aumentação por difusão é mais benéfica quando há poucas amostras por classe, quando as classes são visualmente distintas, quando o modelo de difusão foi bem treinado, e quando a diversidade das gerações complementa o dataset original. Por outro lado, pode não ser benéfica quando há muitas amostras por classe, quando as gerações introduzem artefatos, quando o modelo de difusão não captura bem as características, ou quando a aumentação tradicional já é suficiente.

\section{Eficiência do DGAE vs CNN-JEPA}

\subsection{Comparação de Eficiência}

A comparação entre DGAE e CNN-JEPA revela trade-offs interessantes entre compactação, memória, tempo de treinamento e qualidade de representações. DGAE oferece máxima compactação com espaço latente 2x menor, conforme reportado por \cite{liu2025dgae}, mas requer modelo de difusão em memória durante treinamento, aumentando uso de memória. O tempo de treinamento é maior devido ao guiamento por difusão, mas as representações resultantes são mais compactas e expressivas.

CNN-JEPA oferece eficiência em termos de memória e tempo de treinamento, não requerendo modelos externos, mas o espaço de features é comparável a outras técnicas (512 dimensões), não oferecendo a mesma compactação que DGAE. O treinamento é mais rápido e simples, mas pode requerer mais épocas para convergência.

\subsection{Trade-offs e Aplicações}

A escolha entre DGAE e CNN-JEPA depende dos requisitos específicos da aplicação. DGAE é mais adequado quando se busca máxima compactação, quando já se tem modelos de difusão disponíveis, quando o objetivo é integração com modelos de difusão subsequentes, ou quando recursos computacionais permitem treinamento mais complexo. CNN-JEPA é mais adequado quando se quer evitar dependência de modelos externos, quando se busca simplicidade e eficiência, quando aumentação explícita não é desejável, ou quando recursos computacionais são limitados.

\subsection{Resultados e Implicações}

Os resultados demonstram que ambas as técnicas são competitivas em termos de qualidade de representações para clustering, com DGAE oferecendo vantagem em compactação e CNN-JEPA oferecendo vantagem em simplicidade. Esta observação tem implicações práticas importantes: a escolha de técnica deve ser baseada em requisitos específicos da aplicação, não apenas em métricas de qualidade.

Para aplicações onde compactação é crítica, como sistemas de busca por similaridade ou análise de grandes volumes de dados, DGAE é a escolha mais adequada. Para aplicações onde simplicidade e eficiência são prioritárias, CNN-JEPA oferece melhor balanceamento. Para aplicações onde qualidade máxima é o objetivo principal, SimCLR com aumentação por difusão apresenta melhor desempenho geral.

\subsection{Implicações Práticas}

Os resultados têm implicações práticas importantes para design de pipelines de aprendizado auto-supervisionado. A escolha de técnica deve considerar não apenas métricas de qualidade, mas também requisitos de recursos, disponibilidade de modelos pré-treinados, e objetivos específicos da aplicação. A integração de aumentação por difusão em workflows existentes pode melhorar significativamente qualidade de representações, especialmente em cenários com poucas amostras.

A comparação entre técnicas que dispensam aumentação explícita (CNN-JEPA) e técnicas que a utilizam (SimCLR) demonstra que ambas as abordagens são viáveis, oferecendo flexibilidade no design de sistemas de aprendizado auto-supervisionado. Finalmente, a eficiência do DGAE em termos de compactação abre possibilidades para aplicações onde espaço de armazenamento ou eficiência computacional são críticos.
